\hnsection{线性表示与全局李群}

刚才提到的这些工作没有任何一项真正试图定义和研究全局李群。H. Weyl 首先朝这个方向迈出了步伐。他受到两种理论的启发，这两种理论此前一直独立发展：E. Cartan 关于复半单李代数线性表示的理论，以及 Frobenius 关于有限群线性表示的理论；后者刚刚被 I. Schur 借助 Hurwitz 的一个想法转用于正交群。Hurwitz 已经说明 {[}17{]}，通过把有限群上的平均运算替换为关于不变测度的积分，可以为正交群或酉群构造不变量。他还注意到，将此方法应用于酉群，可以为一般线性群得到不变量，这就是“酉技巧”的第一个例子。1924 年，I. Schur {[}20{]} 用这一过程，通过构造不变的非退化正定 Hermite 型，证明了正交群 \(\mathbf{O}(n)\) 和酉群 \(\mathbf{U}(n)\) 的表示的完全可约性；他由“酉技巧”推出 \(\mathbf{O}(n,\mathbf{C})\) 和 \(\mathbf{SL}(n,\mathbf{C})\) 的全纯表示的完全可约性，建立了 \(\mathbf{O}(n)\) 和 \(\mathbf{U}(n)\) 的特征标正交关系，并确定了 \(\mathbf{O}(n)\) 的特征标。H. Weyl 随即将这种方法推广到复半单李代数 {[}21{]}。给定这样的代数 g，他证明它具有一个“紧实实形式”（这相当于说，它是由 R 上的一个代数 \(g_{0}\) 通过从 R 到 C 的标量扩张得到的，而其伴随群 \(G_{0}\) 是紧的）。此外，他证明 \(G_{0}\) 的基本群是有限的，因而 \(G_{0}\) 的万有覆盖‡是紧的。他通过适当改造 Schur 的方法，推出 g 的表示的完全可约性，并且还用一种全局方法确定了 g 的表示的特征标。在给 I. Schur 的一封信中（Sitzungsber, Berlin, 1924, pp.~338--343），H. Weyl 概述了 Cartan 的、Schur 尚不知道的结果（cf.~{[}20{]}，p.~299，页底注释），并比较了两种方法：Cartan 的方法给出了具有李代数 g 的单连通群的所有全纯表示；对于正交群，这样得到的是一个二重覆盖（后来称为自旋群）的表示，而 Schur 未能得到它；另一方面，Schur 的方法的优点在于显示了完全可约性并给出了特征标的显式形式。

在 H. Weyl 的工作之后，E. Cartan 在关于对称空间和李群的研究中采取了明确的全局观点。这一观点构成了他 1930 年对“有限连续”群理论的论述 ({[}12{]}, vol.~I₂, pp.~1165--1225) 的基础。特别是，这里给出了全局版本的第三基本定理（存在具有给定李代数的李群）的第一个证明；Cartan 还证明，实李群的每个闭子群都是李群（第三章，§ 8，第 2 号，定理 2），从而推广了 J. von Neumann 关于线性群的闭子群的一个结果 {[}23{]}。在这篇回忆录中，von Neumann 还证明了半单群的每个连续表示都是实解析的。

在这些工作之后，“经典”意义下的李群理论（即在 R 或 C 上的有限维理论）在其基本方面大体已经完成。Pontrjagin 在其关于拓扑群的著作中给出了第一个详细论述 {[}36{]}；在那里，他采用了与 Lie 很接近的观点，但仔细区分了局部与全局。随后是 Chevalley 的著作 {[}38{]}，其中还首次系统讨论了解析流形理论和外微分演算；Lie 的“无穷小变换”在那里作为向量场出现，李群 G 的李代数被认同为 G 上左不变向量场的空间。他不讨论“群芽”和“变换群”。

